\section*{Bab \ref{ch:g11:deriv} --- Turunan: Perkenalan Pertama}

\begin{solution}{exo:g11:deriv:1}
Untuk $f(x) = x^2$ di $a = 2$:
\[
\frac{(2+h)^2 - 4}{h} = \frac{4h + h^2}{h} = 4 + h
\xrightarrow[h\to0]{} 4,
\]
jadi $f'(2) = 4$. Untuk $g(x) = x^2 + 3x$ di $a = 1$: $g(1) = 4$ dan
\[
\frac{(1+h)^2 + 3(1+h) - 4}{h} = \frac{5h + h^2}{h} = 5 + h
\xrightarrow[h\to0]{} 5,
\]
jadi $g'(1) = 5$.
\end{solution}

\begin{solution}{exo:g11:deriv:2}
$f'(x) = 12x^2 - 4x + 1$.

$g'(x) = \frac{5}{2\sqrt x} - \frac{3}{x^2}$ (pada $\intoo{0}{+\infty}$).

$h'(x) = 2(x^2 - 3) + (2x+1)(2x) = 2x^2 - 6 + 4x^2 + 2x
= 6x^2 + 2x - 6$ (aturan hasil kali).
\end{solution}

\begin{solution}{exo:g11:deriv:3}
$f(x) = x^2 - 3x$, $a = 2$: $f(2) = -2$, $f'(x) = 2x - 3$ jadi $f'(2) = 1$;
\omterm{def:g11:deriv:tangent}{garis singgungnya} $y = -2 + (x - 2) = x - 4$.

$f(x) = \frac1x$, $a = 1$: $f(1) = 1$, $f'(1) = -1$; \omterm{def:g11:deriv:tangent}{garis singgungnya}
$y = 1 - (x - 1) = 2 - x$.

$f(x) = \sqrt x$, $a = 4$: $f(4) = 2$, $f'(4) = \frac{1}{2\sqrt4} =
\frac14$; \omterm{def:g11:deriv:tangent}{garis singgungnya} $y = 2 + \frac14(x - 4) = \frac{x}{4} + 1$.
\end{solution}

\begin{solution}{exo:g11:deriv:4}
$f'(x) = \frac{(x-1) - (x+2)}{(x-1)^2} = \frac{-3}{(x-1)^2}$.

$g'(x) = \frac{2x(x^2+1) - x^2 \times 2x}{(x^2+1)^2}
= \frac{2x}{(x^2+1)^2}$.

$h'(x) = -\frac{2x+1}{(x^2+x+1)^2}$ (aturan hasil bagi dengan $u = 1$).
\end{solution}

\begin{solution}{exo:g11:deriv:5}
$f'(x) = 3x^2 - 12x + 9 = 3(x^2 - 4x + 3) = 3(x-1)(x-3)$: positif
di luar $\intcc{1}{3}$, negatif di dalamnya. Jadi $f$ naik pada
$\intoc{-\infty}{1}$, turun pada $\intcc{1}{3}$, naik pada
$\intco{3}{+\infty}$, dengan \omterm{def:g10:functions:extrema}{maksimum} lokal $f(1) = 1 - 6 + 9 - 2 = 2$ dan
\omterm{def:g10:functions:extrema}{minimum} lokal $f(3) = 27 - 54 + 27 - 2 = -2$.
\end{solution}

\begin{solution}{exo:g11:deriv:6}
Pada $\intoo{-1}{+\infty}$:
\[
f'(x) = \frac{2x(x+1) - (x^2+3)}{(x+1)^2}
= \frac{x^2 + 2x - 3}{(x+1)^2}
= \frac{(x+3)(x-1)}{(x+1)^2}.
\]
Untuk $x > -1$, baik $x + 3$ maupun $(x+1)^2$ bernilai positif, sehingga
$f'(x)$ bertanda sama dengan $x - 1$: jadi $f$ turun pada $\intoc{-1}{1}$ dan
naik pada $\intco{1}{+\infty}$, dengan \omterm{def:g10:functions:extrema}{minimum} $f(1) = \frac{4}{2} = 2$.
\end{solution}

\begin{solution}{exo:g11:deriv:7}
$f'(x) = 2x - 4$.

\emph{1.} \omterm{def:g11:deriv:tangent}{Garis singgung} mendatar: $f'(x) = 0$ di $x = 2$; titiknya adalah
$(2, f(2)) = (2, 1)$ (titik puncak \omterm{def:g11:quad:parabola}{parabolanya}).

\emph{2.} Sejajar dengan $y = 2x + 1$ berarti bergradien $2$: $2x - 4 = 2$,
jadi $x = 3$, yang memberi titik $(3, f(3)) = (3, 2)$.
\end{solution}

\begin{solution}{exo:g11:deriv:8}
Dengan lebar $x$ (dua sisi berpagar) dan panjang $60 - 2x$:
$A(x) = x(60 - 2x) = 60x - 2x^2$ pada $\intoo{0}{30}$. Lalu
$A'(x) = 60 - 4x$, positif sebelum $x = 15$ dan negatif sesudahnya: luasnya
\omterm{def:g10:functions:extrema}{maksimum} untuk $x = 15$, yang memberi kandang $15 \times 30$ m berluas
$450$ m$^2$. Rumus titik puncak $\alpha = -\frac{60}{2 \times (-2)} = 15$
pada \cref{ch:g11:quad} memberi jawaban yang sama tanpa kalkulus.
\end{solution}

\begin{solution}{exo:g11:deriv:9}
Misalkan $f(x) = \frac{x+1}{2} - \sqrt x$ pada $\intoo{0}{+\infty}$. Maka
\[
f'(x) = \frac12 - \frac{1}{2\sqrt x} = \frac{\sqrt x - 1}{2\sqrt x},
\]
yang bertanda sama dengan $\sqrt x - 1$: negatif pada $\intoo{0}{1}$, positif
pada $\intoo{1}{+\infty}$. Jadi $f$ turun lalu naik, dengan \omterm{def:g10:functions:extrema}{minimum}
$f(1) = 1 - 1 = 0$. Karena itu $f(x) \geq 0$ untuk semua $x > 0$, artinya
$\sqrt x \leq \frac{x+1}{2}$, dengan kesamaan hanya di $x = 1$.
\end{solution}

\begin{solution}{exo:g11:deriv:10}
Dari $\pi r^2 h = 500$ diperoleh $h = \frac{500}{\pi r^2}$, sehingga
\[
S(r) = 2\pi r^2 + 2\pi r \cdot \frac{500}{\pi r^2}
= 2\pi r^2 + \frac{1000}{r},
\qquad
S'(r) = 4\pi r - \frac{1000}{r^2}.
\]
$S'(r) = 0$ ketika $4\pi r^3 = 1000$, yaitu $r^3 = \frac{250}{\pi}$, jadi
$r = \sqrt[3]{250/\pi}$ ($\approx 4.3$ cm). Karena $S'(r) =
\frac{4\pi r^3 - 1000}{r^2}$ bernilai negatif sebelum nilai itu dan positif
sesudahnya, nilai itu adalah \omterm{def:g10:functions:extrema}{minimumnya}.
\end{solution}

\begin{solution}{exo:g11:deriv:11}
\omterm{def:g11:deriv:tangent}{Garis singgung} $y = x^2$ di titik yang berabsis $a$ adalah
$y = a^2 + 2a(x - a) = 2ax - a^2$. Garis itu melalui $P(1, -3)$ ketika
$-3 = 2a - a^2$, yaitu $a^2 - 2a - 3 = 0$, yang \omterm{def:g11:quad:discriminant}{akarnya} $a = 3$ dan
$a = -1$ ($\Delta = 16$). Jadi ada \emph{dua} \omterm{def:g11:deriv:tangent}{garis singgung} yang melalui
$P$:
\[
y = 6x - 9 \quad (a = 3),
\qquad
y = -2x - 1 \quad (a = -1).
\]
\end{solution}

\begin{solution}{exo:g11:deriv:12}
Dari \cref{ex:g11:deriv:study}, $f$ naik sampai \omterm{def:g10:functions:extrema}{maksimum} lokal
$f(-1) = 3$, turun sampai \omterm{def:g10:functions:extrema}{minimum} lokal $f(1) = -1$, lalu naik
lagi; untuk $x$ positif besar (masing-masing negatif besar), $x^3$ menguasai
dan $f$ mengambil nilai yang sebesar-besarnya (masing-masing sekecil-kecilnya).
Dengan membaca garis mendatar $y = k$ menyeberangi \omterm{def:g10:functions:table}{tabel variasinya}:
\begin{itemize}
  \item $k > 3$ atau $k < -1$: garisnya bertemu kurvanya sekali
        --- $1$ penyelesaian;
  \item $k = 3$ atau $k = -1$: garisnya menyentuh sebuah titik balik dan
        memotong satu cabang lain --- $2$ penyelesaian;
  \item $-1 < k < 3$: garisnya memotong ketiga cabangnya --- $3$
        penyelesaian.
\end{itemize}
\end{solution}

\begin{solution}{pb:g11:deriv:1}
\textbf{1.} $f'(x) = 6x - 5$; $g'(x) = 3x^2 - 12$. Dari definisinya:
$\frac{(1 + h)^2 - 1}{h} = 2 + h \to 2$: jadi \omterm{def:g11:deriv:derivative}{turunan}
$x^2$ di $1$ adalah $2$.

\textbf{2.} $g(2) = 8 - 24 = -16$ dan $g'(2) = 12 - 12 = 0$:
\omterm{def:g11:deriv:tangent}{garis singgungnya} adalah garis mendatar $y = -16$. \omterm{def:g11:deriv:tangent}{Garis singgung}
mendatar mengisyaratkan sebuah titik kritis --- di sini sebuah \omterm{def:g10:functions:extrema}{minimum} lokal,
seperti dibenarkan \omterm{def:g10:functions:table}{tabel variasi} pada pertanyaan 5.

\textbf{3.} \omterm{def:g10:reffunc:affine}{Gradiennya} $2a$ di $(a, a^2)$:
$y = a^2 + 2a(x - a) = 2ax - a^2$.

\textbf{4.} Di $x = 0$: $y = -a^2$: \omterm{def:g11:deriv:tangent}{garis singgungnya} memotong
sumbunya sedalam $a^2$ di bawah titik asal --- persis sejauh
\emph{di bawah} nol seperti $P$ berada di atasnya. Resepnya: untuk menggambar
\omterm{def:g11:deriv:tangent}{garis singgung} di titik yang tingginya $h$, hubungkan titik itu dengan titik
pada sumbu sedalam $h$ di bawah titik asal. (Tanpa perlu \omterm{def:g11:deriv:derivative}{turunan} di
meja gambar.)

\textbf{5.} $g'(x) = 3(x^2 - 4)$: positif pada
$\intoo{-\infty}{-2}$, negatif pada $\intoo{-2}{2}$, dan positif
sesudahnya. Jadi $g$ naik sampai \omterm{def:g10:functions:extrema}{maksimum} lokal $g(-2) = 16$, jatuh
sampai \omterm{def:g10:functions:extrema}{minimum} lokal $g(2) = -16$, lalu naik lagi.

\textbf{6.} $F\left(0, \frac14\right)$ dan $T(0, -a^2)$:
$FT = \frac14 + a^2$.

\textbf{7.} $FP = a^2 + \frac14 = FT$: sama kaki di $F$, sehingga
kedua sudut alasnya sama: $\widehat{FTP} = \widehat{FPT}$.

\textbf{8.} Sinar datang di $P$ tegak, jadi sejajar
dengan garis $(TF)$ (yaitu sumbu $y$). \omterm{def:g11:deriv:tangent}{Garis singgung} $(TP)$ memotong
kedua sejajar itu, sehingga sudut antara sinar dan \omterm{def:g11:deriv:tangent}{garis singgungnya}
sama dengan $\widehat{FTP}$ (sudut berseberangan) --- yang sama dengan
$\widehat{FPT}$, yaitu sudut antara $[PF]$ dan \omterm{def:g11:deriv:tangent}{garis singgungnya}.
Jadi sudutnya sama besar pada kedua sisi \omterm{def:g11:deriv:tangent}{garis singgung} di $P$.

\textbf{9.} Menurut hukum pemantulan, sinar yang tiba secara
tegak di $P$ pergi sepanjang garis yang membentuk sudut sama besar
di sisi lainnya --- menurut pertanyaan 8, sepanjang $[PF]$: yaitu melalui
fokusnya, berapa pun $a$. Sebaliknya, cahaya yang lahir di $F$ meninggalkan
setiap titik cerminnya sejajar dengan sumbunya. Antena mengumpulkan,
lampu sorot memancarkan: terbukti.

\textbf{10.} $a = 1$: \omterm{def:g11:deriv:tangent}{garis singgungnya} $y = 2x - 1$, $T(0, -1)$;
$FP = \sqrt{1 + \left(1 - \frac14\right)^2} =
\sqrt{\frac{25}{16}} = \frac54$; $FT = \frac14 + 1 = \frac54$:
sama kaki, seperti dijanjikan.

\textbf{11.} \omterm{def:g11:deriv:tangent}{Garis singgung} di $x_n$ adalah
$y = f(x_n) + f'(x_n)(x - x_n)$; garis itu memotong $y = 0$ di
$x = x_n - \frac{f(x_n)}{f'(x_n)}$.

\textbf{12.} $x_1 = 1 - \frac{-1}{2} = \frac32$;
$x_2 = \frac32 - \frac{1/4}{3} = \frac{17}{12}$;
$x_3 = \frac{577}{408}$: yaitu \omterm{def:g10:numbers:approx}{hampiran} Heron bagi $\sqrt2$,
alias mesin \omterm{pb:g10:functions:1}{titik tetap} pada \cref{pb:g10:functions:1}.
Identitasnya: $x - \frac{x^2 - 2}{2x} = \frac{2x^2 - x^2 + 2}{2x}
= \frac{x^2 + 2}{2x} = \frac12\left(x + \frac2x\right)$ ---
resep Heron \emph{adalah} metode Newton yang diterapkan pada
$x^2 - 2$, dua ribu tahun lebih awal.

\textbf{13.} $x_1 = 5 - \frac{125 - 100}{75} = \frac{14}{3}
\approx 4.6667$; $x_2 \approx 4.6667 - \frac{1.63}{65.33}
\approx 4.6417$. Kalkulator: $\sqrt[3]{100} = 4.6416\ldots$ ---
empat angka desimal yang benar dalam dua langkah.

\textbf{14.} $f'(2) = 0$: \omterm{def:g11:deriv:tangent}{garis singgung} di titik awalnya
mendatar (pertanyaan 2) dan tidak pernah memotong sumbunya --- jadi
rumusnya membagi dengan nol. Peringatannya: mulailah Newton jauh dari
titik kritis (dan, lebih umum, cukup dekat ke \omterm{def:g11:quad:discriminant}{akar} yang
kamu inginkan).

\textbf{15.} Metode bagi dua memperoleh satu angka biner per langkah;
Newton kira-kira \emph{melipatduakan} banyaknya angka yang benar setiap
langkah (di sini $1 \to 3 \to 6$). Ketika setiap langkahnya mahal --- jutaan
parameter, navigasi waktu nyata --- melipatduakan mengalahkan
merayap, dan itulah sebabnya \omterm{def:g11:deriv:tangent}{garis singgung} berjalan di dalam begitu banyak
kepingan silikon.

\textbf{16.} $d^2 = 900 - w^2$, sehingga
$S(w) = w(900 - w^2) = 900w - w^3$, untuk $0 < w < 30$.

\textbf{17.} $S'(w) = 900 - 3w^2 = 0$ di
$w = \sqrt{300} = 10\sqrt3 \approx 17.3$ cm ($S' > 0$ sebelumnya,
$< 0$ sesudahnya: jadi sebuah \omterm{def:g10:functions:extrema}{maksimum}). Lalu
$d = \sqrt{600} = 10\sqrt6 \approx 24.5$ cm.

\textbf{18.} $\frac{d^2}{w^2} = \frac{600}{300} = 2$, jadi
$d = w\sqrt2$: tingginya adalah lebar kali $\sqrt2$ --- itulah
rasio tukang kayu (dan lukisan membagi diameter menjadi tiga
menghasilkannya dengan jangka, tanpa aljabar di lokasi).

\textbf{19.} Balok persegi: $w = d = \sqrt{450} \approx 21.2$
cm, kekakuannya $w^3 = 450^{3/2} \approx 9\,546$. Balok optimum:
$S = 10\sqrt3 \times 600 = 6\,000\sqrt3 \approx 10\,392$. Keuntungannya:
sekitar $8.9\,\%$ lebih kaku dari gelondong yang sama --- \omterm{def:g11:deriv:derivative}{turunan} itu
membayar ongkos penggergajian.

\textbf{20.} Geometri: sudut sama besar pada \omterm{def:g11:deriv:tangent}{garis singgung} mengubah
sebuah paraboloid menjadi pengumpul cahaya. Numerik: meluncur menuruni
\omterm{def:g11:deriv:tangent}{garis singgung} menyelesaikan \omterm{def:g10:algebra:equation}{persamaan} dengan laju penggandaan angka.
Pengoptimuman: \omterm{def:g11:deriv:derivative}{turunan} yang lenyap menemukan persegi panjang terkuat dari
sebatang gelondong. Satu objek, tiga keahlian --- dan kecembungan tahun
berikutnya akan memberitahukan, sebagai tambahan, di sisi mana dari setiap
\omterm{def:g11:deriv:tangent}{garis singgung} kurvanya berada.

\end{solution}
